The main difference between a classical BM (CBM)  and a quantum BM (QBM), is that the latter employs nodes that are defined by Pauli matrices rather than binary units, and therefore constructs a network defined by a parameterized Hamiltonian:
\begin{equation}
    \mathcal{H}(\Theta) = -\sum_{i} b_i \sigma^z_i - \sum_{\{i, j\}} w_{ij} \sigma^z_i \sigma^z_j.
\end{equation}
In this case, the \ac{QBM} is defined as the Gibbs state
\begin{equation}
    \rho(\Theta, \mathcal{H}) = \frac{\exp{-\beta \mathcal{H}(\Theta)}}{Z(\Theta)},
\end{equation}
where
\begin{equation}
    \cZ(\beta, \cH) = \Tr{e^{-\beta \cH(\vb, \vw)}}.
\end{equation}

The cost function is thus the generalized Kullback-Leibler divergence for density matrices, which corresponds to the quantum relative entropy:
\begin{equation}
    \cS(\sigma \| \rho(\beta, \cH)) = \Tr{\sigma \ln \sigma} - \Tr{\sigma \ln \rho(\beta, \cH)},
\end{equation}
where $\sigma$ is the target density matrix representing the data embedded into a mixed state.

The goal of \ac{QBM} training is to find a set of weights and biases
of the Hamiltonian, such that the Gibbs state approximates the target
density matrix. There are various methods of tackling the training of
\acp{QBM}~\cite{Gacon2022, Zoufal2021, Coopmans2023, Kalis2023,
  Huijgen2023}, and an analogous problem to \ac{QBM} training is
Hamiltonian learning~\cite{Anshu2021}.

\section{Variational Gibbs State Preparation} \label{sec:VGSP}

In this section we will describe how one can prepare a Gibbs state
using a \ac{VQA} with the free energy as a cost function, specifically
discussing the algorithm presented in
Ref.~\cite{Consiglio2023}. Section~\ref{sec:vn_entropy}
and~\ref{sec:free_energy} are intended to serve as a qualitative
description of the estimation of the von Neumann entropy and the free
energy, respectively. Consequently, in Section~\ref{sec:framework}, a
quantitative description of the framework of the \ac{VQA} is carried
out, mathematically linking the von Neumann entropy and energy
expectation, with the estimation of the free energy.

Suppose one chooses a Hamiltonian $\cH$ and inverse temperature
$\beta$. For a general state $\rho$, one can define a generalized
Helmholtz free energy as


\begin{equation}
    \cF(\rho) = \Tr{\cH \rho} - \beta^{-1}\cS(\rho),
    \label{eq:helmholtz_free_energy}
\end{equation}
where the von Neumann entropy $\cS(\rho)$ can be expressed in terms of the eigenvalues, $p_i$, of $\rho$,
\begin{equation}
    \cS(\rho) = -\sum_{i=0}^{d - 1} p_i \ln p_i.
    \label{eq:von_neumann}
\end{equation}
Since the Gibbs state is the unique state that minimizes the free energy of $\cH$~\cite{Matsui1994}, a variational form can be put forward that takes Eq.~\eqref{eq:helmholtz_free_energy} as an objective function, such that
\begin{equation}
    \rho(\beta, \cH) = \underset{\rho}{\arg\min}~\cF(\rho).
\end{equation}
In this case, $p_i = \exp\left(-\beta E_i\right)/\cZ(\beta, \cH)$ is the probability of getting the eigenstate $\ket{E_i}$ from the ensemble $\rho(\beta, \cH)$.

\subsection{Computing the von Neumann Entropy}~\label{sec:vn_entropy}

The difficulty in measuring the von Neumann entropy, defined by
Eq.~\eqref{eq:von_neumann}, of a quantum state on a \ac{NISQ} device
is typically the challenging part of variational Gibbs state
preparation algorithms, since $\cS(\rho)$ is not an observable. To
exactly compute the von Neumann entropy on a quantum device, one would
have to perform full state tomography on the Gibbs state, which has a
time complexity of $\cO(3^n)$ and a space complexity of $\cO(4^n)$. As
a consequence, different works have employed various techniques to
compute an approximation for the von Neumann entropy, such as using: a
Fourier series approximation to the von Neumann
entropy~\cite{Chowdhury2020}; a thermal multi-scale entanglement
renormalization ansatz~\cite{Sewell2022}; sums of R\'enyi
entropies~\cite{Foldager2022}; among others. On the other hand,
Ref.~\cite{Consiglio2023} devised a method of computing the von
Neumann exactly, via post-processing of measurement results acting on
ancillary register.

When preparing an $n$-qubit state, the unitary gates used in quantum
computers ensure that the final quantum state of the entire register,
starting from the input state $\ket{0}^{\otimes n}$, remains pure. As
a result, in order to prepare an $n$-qubit Gibbs state on the system
register, an $m \leq n$-qubit ancillary register is required. For
example, in the case of the infinite-temperature Gibbs state, which is
a fully mixed state, $m = n$ qubits are needed in the ancillary
register to achieve maximal von Neumann entropy. To accurately
evaluate the von Neumann entropy without any approximations, the
complete Boltzmann distribution is prepared on the ancillary register,
thus, $m = n$ is set regardless of the temperature.

\subsection{Computing the Free Energy}~\label{sec:free_energy}

Following the prescription of Ref.~\cite{Consiglio2023}, we shall
denote the ancillary register as $A$, while the preparation of the
Gibbs state will be carried out on the system register $S$. The
purpose of the \ac{VQA} is to effectively create the Boltzmann
distribution on $A$, which is then imposed on $S$, via intermediary
\textsc{CNOT} gates, to generate a diagonal mixed state in the
computational basis. In the ancillary register we can choose a unitary
ansatz capable of preparing such a probability distribution. Thus, the
ancillary qubits are responsible for classically mixing in the
probabilities of the thermal state, while also being able to access
these probabilities via measurements in the computational basis. On
the other hand, the system register will host the preparation of the
Gibbs state, as well as the measurement of the expectation value of
our desired Hamiltonian during the optimization routine.

The specific design of the \ac{PQC} allows classical post-processing
of simple measurement results, carried out on ancillary qubits in the
computational basis, to determine the von Neumann entropy. A
diagrammatic representation of the structure of the \ac{PQC} is shown
in Fig.~\ref{fig:gibbs_circuit}. Note that while the \ac{PQC} of the
algorithm has to have a particular structure --- a unitary acting on
the ancillae and a unitary acting on the system, connected by
intermediary \textsc{CNOT} gates --- it is not dependent on the choice
of Hamiltonian $\cH$, inverse temperature $\beta$, or the variational
ans\"atze, $U_A$ and $U_S$, employed within. This is in addition to
enjoying a sub-exponential scaling in the number of shots needed to
precisely compute the Boltzmann probabilities. A qualitative analysis
of a power-law scaling is presented in
Section~\ref{sec:error_analysis}.

\subsection{Modular Structure of the PQC}~\label{sec:framework}

The \ac{PQC}, as shown in Fig.~\ref{fig:gibbs_circuit} for the
\ac{VQA}, is composed of a unitary gate $U_A$ acting on the ancillary
qubits, and a unitary gate $U_S$ acting on the system qubits, with
\textsc{CNOT} gates in between. Note that the circuit notation we are
using here means that there are $n$ qubits for both the system and the
ancillae, as well as $n$ \textsc{CNOT} gates that act in parallel, and
are denoted as

\begin{equation}
    \textsc{CNOT}_{AS} \equiv \bigotimes\limits_{i=0}^{n - 1}\textsc{CNOT}_{A_i S_i}.
\end{equation}

The parametrized unitary $U_A$ acting on the ancillae, followed by
\textsc{CNOT} gates between the ancillary and system qubits, is
responsible for preparing a probability distribution on the
system. The parametrized unitary $U_S$ is then applied on the system
qubits to transform the computational basis states into the
eigenstates of the Hamiltonian.

\begin{figure}[t]
    \centering
    \includegraphics[width=0.5\textwidth]{pqc.pdf}
    \caption{\ac{PQC} for Gibbs state preparation, with systems $A$ and $S$ each carrying $n$ qubits. \textsc{CNOT} gates act between each qubit $A_i$ and corresponding $S_i$.}
    \label{fig:gibbs_circuit}
\end{figure}

Throughout the paper we will describe unitary operations and density
matrices in the computational basis: $\ket{0} \equiv
\left(1~0\right)^\top$ and $\ket{1} \equiv \left(0~1\right)^\top$,
which is spanned by one qubit (and products of the basis states in the
case of multiple qubits). We will also denote the $n$-fold tensor
product of a state $\ket{\psi}$, of an $n$-qubit register $R$, as
$\ket{\psi}_R^{\otimes n} \equiv \bigotimes_{i = 0}^{n - 1}
\ket{\psi}_i$, where $i \in R$. Denote a general unitary gate of
dimension $d = 2^n$, as $U_A = (u_{i, j})_{0\leq i, j \leq
  d-1}$. Starting with the initial state of the $2n$-qubit register,
$\ket{0}_{AS}^{\otimes 2n}$, we apply the unitary gate $U_A$ on the
ancillae to get a quantum state $\ket{\psi}_A$, such that

\begin{equation}
    (U_A \otimes I_S) \ket{0}_{AS}^{\otimes 2n} = \ket{\psi}_{A} \otimes \ket{0}_{S}^{\otimes n},
    \label{eq:U_A}
\end{equation}

where $\ket{\psi}_A = \sum_{i=0}^{d - 1} u_{i,0} \ket{i}_A$ and $I_S$
is the identity acting on the system. Since Eq.~\eqref{eq:U_A} is
applied to the all-zero state in the ancillary register, then the
operation serves to extract the first column of the unitary operator
$U_A$. The next step is to prepare a classical probability mixture on
the system qubits, which can be done by applying \textsc{CNOT} gates
between each ancilla and system qubit. This results in a state

\begin{align}
    \textsc{CNOT}_{AS} \left( \ket{\psi}_{A} \otimes \ket{0}_{S}^{\otimes n} \right) = \sum_{i=0}^{d - 1} u_{i,0} \ket{i}_A \otimes \ket{i}_S.
    \label{eq:pre_tfd}
\end{align}
By then tracing out the ancillary qubits, we arrive at
\begin{align}
    \PTr{A}{\left( \sum_{i=0}^{d-1} u_{i,0}\ket{i}_A \otimes \ket{i}_S \right) \left( \sum_{j=0}^{d-1} u_{j,0}^* \bra{j}_A \otimes \bra{j}_S \right)} = &\sum_{i,j=0}^{d-1} u_{i,0}u_{j,0}^* \braket{i}{j} \ketbra{i}{j}_S \nonumber \\ = \
    &\sum_{i=0}^{d-1} |u_{i,0}|^2 \ketbra{i}{i}_S,
    \label{eq:diag}
\end{align}

ending up with a diagonal mixed state on the system, with
probabilities given directly by the absolute square of the entries of
the first column of $U_A$, that is, $p_i = |u_{i,0}|^2$. If the system
qubits were traced out instead, we would end up with the same diagonal
mixed state,

\begin{align}
    \PTr{S}{\left( \sum_{i=0}^{d-1} u_{i,0}\ket{i}_A \otimes \ket{i}_S \right) \left( \sum_{j=0}^{d-1} u_{j,0}^* \bra{j}_A \otimes \bra{j}_S \right)} = &\sum_{i,j=0}^{d-1} u_{i,0}u_{j,0}^* \braket{i}{j} \ketbra{i}{j}_A \nonumber \\ = \
    &\sum_{i=0}^{d-1} |u_{i,0}|^2 \ketbra{i}{i}_A.
\end{align}

This implies that by measuring in the computational basis of the
ancillary qubits, we can determine the probabilities $p_i$, which can
then be post-processed to determine the von Neumann entropy $\cS$ of
the state $\rho$ via Eq.~\eqref{eq:von_neumann} (since the entropy of
$A$ is the same as that of $S$). As a result, since $U_A$ only serves
to create a probability distribution from the entries of the first
column, we can do away with a parametrized orthogonal (real unitary)
operator, thus requiring less gates and parameters for the ancillary
ansatz.

The unitary gate $U_S$ then serves to transform the computational
basis states of the system qubits to the eigenstates of the Gibbs
state, once it is optimized by the \ac{VQA}, such that


\begin{align}
    \rho = U_S \left( \sum_{i=0}^{d-1} |u_{i,0}|^2 \ketbra{i}{i}_S \right) U_S^\dagger = \sum_{i=0}^{d - 1} p_i \ketbra{\psi_i},
\end{align}

where the expectation value $\Tr{\cH \rho}$ of the Hamiltonian can be
measured. Ideally, at the end of the optimization procedure, $p_i =
\exp\left(-\beta E_i\right)/\cZ(\beta, \cH)$ and $\ket{\psi_i} =
\ket{E_i}$, so that we get

\begin{equation}
    \rho(\beta, \cH) = \sum_{i=0}^{d - 1} \frac{e^{-\beta E_i}}{\cZ(\beta, \cH)} \ketbra{E_i}.
\end{equation}

The \ac{VQA} therefore avoids the entire difficulty of measuring the
von Neumann entropy of a mixed state on a quantum computer, and
instead transfers the task of post-processing measurement results to
the classical computer, which is much more tractable.

\subsection{Objective Function}

Finally, we can define the objective function of our \ac{VQA} to
minimize the free energy~\eqref{eq:helmholtz_free_energy}, via our
constructed \ac{PQC}, to obtain the Gibbs state

\begin{align}
    \rho(\beta, \cH) &= \underset{\bm{\theta}, \bm{\varphi}}{\arg\min} \ \cF\left(\rho\left(\bm{\theta}, \bm{\varphi}\right)\right) \nonumber \\ &= \underset{\bm{\theta}, \bm{\varphi}}{\arg\min} \left( \Tr{\cH \rho_S(\bm{\theta}, \bm{\varphi})} - \beta^{-1}\cS\left(\rho_A(\bm{\theta})\right) \right).
    \label{eq:gibbs_state}
\end{align}

It is noteworthy to mention that while the energy expectation depends
on both sets of angles $\bm{\theta}$ (as $U_A$ is responsible for
parameterizing the Boltzmann distribution) and $\bm{\varphi}$ (as
$U_S$ is responsible for parameterizing the eigenstates of the Gibbs
state), the calculation of the von Neumann entropy only depends on
$\bm{\theta}$.

Furthermore, once we obtain the optimal parameters $\bm{\theta^*}$,
$\bm{\varphi^*}$, preparing the Gibbs state $\rho(\beta, \cH)$ on the
system qubits $S$, one can place the same unitary $U_S$ with optimal
parameters $\bm{\varphi^*}$ on the ancillary qubits to prepare the
\ac{TFD} state on the quantum computer, as shown in
Fig.~\ref{fig:tfd_circuit}. Notice that this is equivalent to
Eq.~\eqref{eq:pre_tfd} followed by applying the optimized $U_S$ on
both registers. Tracing out either the ancilla or system register
yields the same Gibbs state on the other register.

\begin{figure}[t]
    \centering
    \includegraphics[width=0.6\textwidth]{tfd.pdf}
    \caption{Optimal \ac{PQC} for \ac{TFD} state preparation, with systems $A$ and $S$ each carrying $n$ qubits. \textsc{CNOT} gates act between each qubit $A_i$ and corresponding $S_i$.}
    \label{fig:tfd_circuit}
\end{figure}

\section{Preparing Gibbs states of the \textit{XY} model} \label{sec:results}

In this Section we assess the performance of the \ac{VQA} for Gibbs
state preparation of an $XY$ model. The $XY$ model~\cite{Lieb1961} is
defined as

\begin{equation}
    \cH = -\sum_{i=1}^n \left( \frac{1 + \gamma}{2} \sigma_i^x \sigma_{i+1}^x + \frac{1 - \gamma}{2} \sigma_i^y \sigma_{i+1}^y \right) - h \sum_{i=1}^n \sigma_i^z.
    \label{eq_XY}
\end{equation}

Here we only report one relevant property for implementing a
problem-inspired ansatz for $U_S$. The Hamiltonian in
Eq.~\eqref{eq_XY} commutes with the parity operator
$\cP=\prod_{i=0}^{n - 1} \sigma^z_i$. As a consequence, the
eigenstates of $\cH$ have definite parity, and so will the eigenstates
of $\rho_{\beta}$.

To assess the performance of the \ac{VQA}, we utilize the
Uhlmann-Josza fidelity as a figure of merit~\cite{Uhlmann2011},
defined as $F\left(\rho,\sigma\right) =
\left(\Tr{\sqrt{\sqrt{\rho}\sigma\sqrt{\rho}}}\right)^2$. This
fidelity measure quantifies the ``closeness'' between the prepared
state and the target Gibbs state, making it a commonly-employed metric
for distinguishability. However, alternative measures exist, each with
its own interpretation. For instance, the trace
distance~\cite{Nielsen2010} can be used, which guarantees that if its
value between two states is bounded by $\epsilon$, the expectation
values computed on the effectively prepared state will differ from
those of the true Gibbs state by at most an amount proportional to
$\epsilon$\cite{Holmes2022}. Another option is the relative
entropy\cite{Nielsen2010}, which characterizes the distinguishability
between the two states as the surprise that arises when an event
occurs that is not possible with the true Gibbs
state~\cite{Vedral2002}.

We use an alternating, entangling brick-wall \ac{PQC} for the unitary
$U_A$, composed of parametrized $R_Y(\theta_i)$ gates, and
\textsc{CNOT}s as the entangling gates. This ansatz is hardware
efficient and is sufficient to produce real amplitudes for preparing
the probability distribution. Note that we require the use of
entangling gates~\cite{Consiglio2023}, as otherwise we will not be
able to prepare any arbitrary probability distribution, including the
Boltzmann distribution of the $XY$ model.

For the unitary $U_S$, We employ a brick-wall structure solely using
parity-preserving gates --- $R_{XY}(\varphi_i)$ followed by
$R_{YX}(\varphi_j)$ gates. If we combine the two gates, which we
denote as $R_P(\varphi_i, \varphi_j)$, we get

\begin{align}
    R_P(\varphi_i, \varphi_j) &= R_{YX}(\varphi_j)\cdot R_{XY}(\varphi_i) \nonumber \\
    &=
    \left(
    \begin{array}{cccc}
     \cos \left(\frac{\varphi_i +\varphi_j }{2}\right) & 0 & 0 & \sin \left(\frac{\varphi_i +\varphi_j }{2}\right) \\
     0 & \cos \left(\frac{\varphi_i -\varphi_j }{2}\right) & -\sin \left(\frac{\varphi_i -\varphi_j }{2}\right) & 0 \\
     0 & \sin \left(\frac{\varphi_i -\varphi_j }{2}\right) & \cos \left(\frac{\varphi_i -\varphi_j }{2}\right) & 0 \\
     -\sin \left(\frac{\varphi_i +\varphi_j }{2}\right) & 0 & 0 & \cos \left(\frac{\varphi_i +\varphi_j }{2}\right) \\
    \end{array}
    \right),
    \label{eq:R_P}
\end{align}

which can be decomposed into two \textsc{CNOT} gates six $\sqrt{X}$
gates, and ten $R_Z$ gates. The decomposed unitary is shown in
Fig.~\ref{fig:R_P}. One layer of the unitary acting on the system
qubits consists of a brick-wall structure, composed of an even-odd
sublayer of $R_P$ gates, followed by an odd-even sublayer of $R_P$
gates. Table~\ref{tab:scaling} shows the scaling of the \ac{VQA}
assuming a closed ladder connectivity, for $n > 2$. An example of a
\ac{PQC}, split into a four-qubit ancilla register, and a four-qubit
system register, is shown in Fig.~\ref{fig:pqc_4qubits}.

\begin{figure}[t]
    \centering
    \includegraphics[width=\textwidth]{R_P.pdf}
    \caption{Decomposed $R_P$ gate in Eq.~\eqref{eq:R_P}.}
    \label{fig:R_P}
\end{figure}

\begin{table}[t]
\caption{Scaling of the \ac{VQA} assuming a closed ladder connectivity, for $n > 2$, where $l_A$ and $l_S$ are the number of ancilla ansatz and system ansatz layers, respectively, and $P$ is 2 when $n$ is even and 3 when $n$ is odd. Circuit depth is calculated on the depth of \textsc{CNOT} gates.}
\centering
\begin{tabular}{|l|l|l|} 
    \hline
    \# of parameters & $n(l_A + 1) + 2nl_S$ & $\cO(n(l_A + l_S))$ \\ 
    \hline
    \# of \textsc{CNOT} gates & $nl_A + 2nl_S + n$ & $\cO(n(l_A + l_S))$ \\
    \hline
    \# of $\sqrt{X}$ gates & $2n(l_A + 1) + 6nl_S$ & $\cO(n(l_A + l_S))$ \\ 
    \hline
    Circuit depth & $P l_A + 2Pl_S + 1$ & $\cO(l_A + l_S)$ \\ 
    \hline
\end{tabular}
\label{tab:scaling}
\end{table}

\begin{figure}[t]
    \centering
    \includegraphics[width=\textwidth]{pqc_4qubits.pdf}
    \caption{Example of an eight-qubit \ac{PQC}, consisting of three ancilla layers acting on a four-qubit register, and three system layers acting on another four-qubit register. Each $R_Y$ gate is parametrized with one parameter $\theta_i$, while each $R_P$ gate has two parameters $\varphi_i$ and $\varphi_j$. The $R_P$ gate is defined in Eq.~\eqref{eq:R_P} with its decomposition shown in Fig.~\ref{fig:R_P}. Note that the intermediary \textsc{CNOT} gates, as well as the \textsc{CNOT} and $R_P$ gates acting on qubits two and three, and on qubits one and four of the ancilla and system, respectively, can be carried out in parallel, respectively.}
    \label{fig:pqc_4qubits}
\end{figure}

\subsection{Statevector Results} \label{sec:statevector_results}

Fig.~\ref{fig:statevector_fidelity} shows the fidelity of the generated mixed state when compared with the exact Gibbs state of the $XY$ model with $h = 0.5$, and $\gamma$ ranging from 0.1 to 0.9 in steps of 0.1, across a range of temperatures for system sizes between two to seven qubits. The \ac{VQA} was carried out using statevector simulations with the \ac{BFGS} optimizer~\cite{Nocedal2006}. We used $n - 1$ layers for both the ancilla ansatz and for the system ansatz, with the scaling highlighted in Table~\ref{tab:scaling_2}. The number of layers was heuristically chosen to satisfy, at most, a polynomial scaling in quantum resources (but linear in depth), while achieving a fidelity higher than 99\%. Furthermore, in order to alleviate the issue of getting stuck in local minima, the optimizer was embedded in a Monte Carlo framework, that is, taking ten random initial positions and carrying out a local optimization from each position --- which we call a `run' --- and finally taking the global minimum to be the minimum over all runs.

\begin{figure}[!ht]
    \centering
    \includegraphics[width=\textwidth]{9fidelity_statevector.pdf}
    \caption{Fidelity $F$, of the obtained state via statevector simulations (using \ac{BFGS}) with the exact Gibbs state, vs inverse temperature $\beta$, for two to seven qubits of the $XY$ model with $\gamma$ between 0.1 and 0.9, and $h = 0.5$. A total of ten runs are made for each point, with the optimal state taken to be the one that maximizes the fidelity.}
    \label{fig:statevector_fidelity}
\end{figure}

\begin{table}[t]
\caption{Scaling of the \ac{VQA} assuming a closed ladder connectivity, for $n > 2$, with $l_A = n - 1$, and $l_S = n - 1$, and $P$ is 2 when $n$ is even and 3 when $n$ is odd. Circuit depth is calculated on the depth of \textsc{CNOT} gates.}
\centering
\begin{tabular}{|l|l|l|} 
    \hline
    \# of parameters & $n(3n - 2)$ & $\cO(n^2)$ \\ 
    \hline
    \# of \textsc{CNOT} gates & $3n^2 - 2n$ & $\cO(n^2)$ \\
    \hline
    \# of $\sqrt{X}$ gates & $2n(4n - 3)$ & $\cO(n^2)$ \\ 
    \hline
    Circuit depth & $3P(n - 1) + 1$ & $\cO(n)$ \\
    \hline
\end{tabular}
\label{tab:scaling_2}
\end{table}

A total of ten runs of \ac{BFGS} per $\beta$ were carried out to verify the reachability of the \ac{PQC}, with Fig.~\ref{fig:statevector_fidelity} showcasing the maximal fidelity achieved for each $\beta$ out of all runs. The results show that, indeed, our \ac{VQA}, is able to reach a very high fidelity $F > 99\%$ for up to seven-qubit Gibbs states of the $XY$ model. In the case of the extremal points, that is $\beta \rightarrow 0$ and $\beta \rightarrow \infty$, the fidelity reaches unity, for all investigated system sizes. For $\beta \rightarrow 0$, the problem reduces to maximizing the von Neumann entropy, where the problem becomes independent of the energy, and therefore, independent of the prepared eigenstates by $U_S$. For $\beta \rightarrow \infty$, the problem reduces to minimizing the energy, which is equivalent to finding the ground-state of the Hamiltonian, and so, independent of $U_A$. This is in contrast with intermediary temperatures $\beta \sim 1$, where the fidelity decreases with the number of qubits. This could be attributed to the fact that at intermediate temperatures, a substantial amount of Boltzmann probabilities are necessary for evaluating the von Neumann entropy with a high precision, and similarly, a large number of eigenstates contribute to preparing the Gibbs state, resulting in a dip in the fidelity at those temperatures.

